\section{Results and scope}\label{inf:sec-introduction}

In the smooth sphere-boundary formulation, the Schiffer conjecture
states that a bounded domain $\Omega\subset\R^n$ with boundary
homeomorphic to $S^{n-1}$ must be a ball if it supports a nonconstant
solution of $\Delta u+\lambda u=0$, for some $\lambda>0$, with
constant boundary value and zero normal derivative
\cite{inf:Wil76,inf:Zal92}. We normalize the frequency to one and
the boundary value to one. We construct domains with these
overdetermined data in unbounded families of both odd and even ambient
dimensions, while retaining strict convexity.

Write $J_\nu$ for the Bessel function of the first kind of order $\nu$.
The real primary crossing orders $\widehat p_m$, for all sufficiently
large indices $m\ge m_0$, satisfy
\[
 J_{\widehat p_m}(2\sqrt{(\widehat p_m+1)(\widehat p_m+2)})=0.
\]
They are specified uniquely by the expansion in
Proposition~\ref{inf:primary-expansion}. The arithmetic construction
uses the fixed irrational number
\[
 \alpha_{\mathrm{ar}}=\frac{\pi}{2(\sqrt3-\pi/3)}
\]
and continued-fraction approximations with denominator one modulo
four. Define $\mathcal P_{\mathrm{bd}}$ to consist of all positive
$p\in\tfrac12\Z$ with $|p-\widehat p_m|<31/(10\widehat p_m)$
for some $m\ge m_0$. The detuning is \eqref{inf:detuning-definition}
at the Neumann zero continued from the nearby crossing.
For every constructed domain the frequency-one
Schiffer data mean a nonconstant classical solution of
\begin{equation}\label{inf:main-equations}
 \Delta u+u=0\quad\hbox{in }\Omega,\qquad
 u=1,\quad\nabla u=0\quad\hbox{on }\partial\Omega.
\end{equation}
\begin{maintheorem}[Bounded detuning]\label{inf:main}\label{inf:theorem-A}
For every sufficiently large $p\in\mathcal P_{\mathrm{bd}}$, there
is an integer $a\in[p/2,3p/4]$ and a bounded strictly convex non-ball
domain $\Omega\subset\R^{2p}$ with real-analytic boundary,
invariant under $O(a)\times O(2p-a)$ and carrying
\eqref{inf:main-equations}. Its detuning satisfies $0<|\theta|\le20$.
There are infinitely many such inputs in each of $\Z$ and $\Z+1/2$.

For every fixed positive integer $L_{\mathrm{ap}}$ and integer
$b_{\mathrm{ap}}$, there is also an unbounded selection
$n\equiv b_{\mathrm{ap}}\pmod{4L_{\mathrm{ap}}}$ with
$|n/2-\widehat p_m|<32L_{\mathrm{ap}}^2/\widehat p_m$ and
$0<|\theta|<256L_{\mathrm{ap}}^2$, whose sufficiently large members
have the same domain conclusion. The threshold may depend on the
fixed modulus and residue. Hence every residue modulo
$L_{\mathrm{ap}}$ contains infinitely many such dimensions.
For even $L_{\mathrm{ap}}$ a residue modulo $L_{\mathrm{ap}}$ permits
only its compatible parity; for odd $L_{\mathrm{ap}}$ each residue
modulo $L_{\mathrm{ap}}$ contains infinitely many dimensions of both parities.

In each of these bounded families, fix
$0<\varepsilon_{\mathrm{sp}}<1/4$ and $0<\eta_{\mathrm{cnt}}<1/3$.
For every sufficiently large member $n=2p$, all but
$O(n^{2/3+\eta_{\mathrm{cnt}}})$ integer splits
$\varepsilon_{\mathrm{sp}}n\le a\le(1/2-\varepsilon_{\mathrm{sp}})n$
give such domains, pairwise non-similar for distinct retained splits.
The constant and threshold may depend on the family and these fixed
parameters. Let $\mathcal N_{\mathrm{sim}}(n)$ be the number of
integer splits $0<a<n/2$ for which the construction on some fixed
compact interior interval gives a domain. Choosing one domain at
each such split gives this many similarity classes, independently
of the choices.
Then along each family's dimensions
\[
 \liminf_{n\to\infty}\mathcal N_{\mathrm{sim}}(n)/n\ge1/2.
\]
The closure uses the cofactor estimate and a fixed centre order
$N=4M+40$, with $M$ fixed after the bounded family and compact
split interval have been chosen and before $N$.
\end{maintheorem}

\begin{maintheorem}[Logarithmic window]\label{inf:theorem-B}
Let
\[
 \mathcal P_{\log}=\{p\in\tfrac12\Z:p>1,\quad
 \exists m\ge m_0:\ |p-\widehat p_m|\le\log p/(8p)\}.
\]
For every sufficiently large $p\in\mathcal P_{\log}$, with $n=2p$, there is an
integer split $0<a<n/2$ and a bounded strictly convex non-ball domain
$\Omega\subset\R^{2p}$ with real-analytic boundary, invariant under
$O(a)\times O(2p-a)$ and carrying \eqref{inf:main-equations}.
For every fixed $0<\varepsilon_{\mathrm{sp}}<1/4$ and
$0<\eta_{\mathrm{cnt}}<1/3$, all but
$O_{\mathrm{sp},\eta_{\mathrm{cnt}}}(n^{2/3+\eta_{\mathrm{cnt}}})$
interior integer splits in the interval of Theorem~\ref{inf:theorem-A}
give pairwise non-similar domains. The same similarity-class lower
limit is at least $1/2$. This family contains sequences of both
parities with $|\theta|\to\infty$. The closure uses mixed
conditioning and the fixed centre order $N=20$.
\end{maintheorem}

\begin{proof}[Proofs of Theorems~\ref{inf:theorem-A} and \ref{inf:theorem-B}]
Corollary~\ref{inf:bounded-basic} proves the first bounded construction.
Lemmas~\ref{inf:many-splits} and \ref{inf:similarity-classes},
Proposition~\ref{inf:bounded-assembly}, and
Theorem~\ref{inf:residue-arithmetic} give all the further assertions
of Theorem~\ref{inf:theorem-A} through the cofactor route.
Proposition~\ref{inf:log-assembly} gives the logarithmic construction
and multiplicity, and Lemma~\ref{inf:log-multiplier} gives unbounded
detuning in both parities. These prove Theorem~\ref{inf:theorem-B}.
\end{proof}

Both assertions concern sufficiently large inputs, with implicit
thresholds and splits. Each fixed-width family and the logarithmic
family have density zero
by Proposition~\ref{inf:arithmetic-density}. The quartic profile and
its Hermite limit are described in Section~\ref{inf:sec-shape}.

The Pompeiu property means that a continuous function on $\R^n$
whose integral over every rigid copy of a domain is zero must
itself be zero. The Pompeiu conjecture states that, among bounded
domains with Lipschitz boundary homeomorphic to $S^{n-1}$,
failure of this property forces the domain to be a ball
\cite{inf:Wil76,inf:Zal92}. The domains in both theorems fail this property
because their indicator transforms vanish on the unit sphere.
A single nonzero plane wave, or its real part, then gives the
required nonzero function. This implication is proved directly
by Green's identity in Proposition~\ref{inf:Pompeiu}.
Consequently the Schiffer and Pompeiu conjectures fail within
the class of convex domains in infinitely many dimensions.

Brown, Schreiber and Taylor developed the spectral-synthesis
approach to the Pompeiu problem and its Fourier--Laplace
characterization \cite{inf:BST73}. Williams related failure of the
Pompeiu property on simply connected Lipschitz domains to the
overdetermined eigenvalue problem \cite{inf:Wil76};
Zalcman's survey gives the classical
bibliography \cite{inf:Zal92}. Recent planar counterexamples were
constructed by Cao-Labora and de Dios Pont \cite{inf:CLdDP26}
and by Colbrook and Stepaniants \cite{inf:CS26}.
Convex counterexamples in the plane, in every dimension from three
to eighteen, and in dimensions twenty and twenty-one were obtained
by computer-assisted proofs in \cite{inf:Guo26}.
The construction below uses a two-block split parameter at fixed
ambient dimension and closes finite-order centres in a graded
Jacobi coefficient space.

\subsection{The construction and its estimates}

The block split supplies a parameter while the ambient dimension
and the ball radial spectrum remain fixed. Write an ambient point as
$(X,Y)\in\R^a\times\R^{2p-a}$. With $n=2p$ and
$r=a/(2p)$, invariant functions depend on radius and
$x=|X|^2/(|X|^2+|Y|^2)$. Their angular basis is Jacobi orthogonal
for the beta measure of parameters $pr,p(1-r)$.
The monic quadratic $\mathsf H_2(x)$ has spherical degree four and is the
shape used to start the construction.
At a continued radial Neumann zero $J_p(R)=0$, define
\[
 \chi=\frac{J_{p+3}(R)}{J'_{p+3}(R)-(p-1)J_{p+3}(R)/R},
 \qquad \theta=4Rp\chi.
\]
For all sufficiently large members of the bounded base family,
$R=2p+3+O(p^{-1})$ and $0<|\theta|<20$.
The estimates are uniform for $0<|\theta|\le20$.

The normalized quartic radius is
$1-\theta \mathsf H_2/(2R^2pg_2)$, where
$g_2=[\mathsf H_2]\mathsf H_2^2\asymp p^{-1}$.
Vary the real split through $\varsigma=(r-1/2)^2\in[1/64,1/16]$.
The exact weighted Dirichlet eigenvalue curves have signed speed
at least $c|\theta|/p$ when their offsets from $R^2$ lie in
$(-|\theta|/(16p),|\theta|/(16p))$.
Only $O(p^{2/3})$ ordered curves can enter this band.
Avoiding their small intervals on the integer split grid gives an
actual invariant Dirichlet gap $c|\theta|p^{-5/3}$ for at least
one admissible $a$.

For each fixed accuracy order $N$, finite-order centre construction
at that split solves the interior
Helmholtz equation exactly and makes both boundary defects
$O_N(\theta^2p^{-N})$ on every prescribed finite complex disc,
with the constant and threshold allowed to depend on that disc.
The recurrence uses only fixed nonzero leading divisors, including
the seed divisor $-1/8$ after its natural scaling.
The centre differs from the quartic radius by
$O_N(|\theta|p^{-4})$ in normalized coordinates, so it inherits
half the Dirichlet gap by domain monotonicity.

The nonlinear argument uses one fixed coefficient space
$\Cc_{p,2}$ for sources and a boundary space $\Ec_1$ with one
angular degree moment. The clamped source subspace $\Gc_p$
contains exactly the sources whose Poisson fields have zero
Dirichlet and normal traces. Positive Jacobi multiplication and a
parameter-raising coupling give a high-shape-degree gain on these
clamped fields. This makes the scalar pullback analytic on the
fixed small-shape chart in $\Gc_p\times\Ec_1$, with a
second-derivative bound $C_Np^6$ on the stated centre neighbourhood.
An exact material change of variables has inverse bounded by
$C_Np^6$ in those same spaces.

The remaining inverse estimate passes through
$\Mmix$, which has radial $L^2$ norms and angular $\ell^1$
coefficients. In a ball window of width $E_{\mathrm{inv}}|\theta|$,
the exact Schur matrix has entries of size $O(|\theta|)$,
exponential angular decay in two matrix norms, and the actual
Hilbert gap. Its logarithmic angular clusters have bounded
\emph{cardinality} $M$, independent of $p$ and of the separation
constant. A determinant and cofactor estimate on each cluster
therefore gives
\[
 \nrm{F_{\mathrm{eff}}^{-1}}_{\ell^1(2^j)}
 \le C|\theta|^{-1}p^{5M/3}.
\]
Radial graph reconstruction and a finite-head/infinite-tail
argument give the full coefficient inverse for the scalar Dirichlet
graph operator $J_D=B_{\Phi_N,D}\KD$, with $B_{\Phi,D}$ defined after
\eqref{inf:scalar-pullback-definition},
\begin{equation}\label{inf:intro-inverse}
 \nrm{J_D^{-1}}_{\Cc_{p,2}\to\Cc_{p,2}}
 \le C|\theta|^{-1}p^{5/2+5M/3}.
\end{equation}

The constants determining $M$ and the inverse exponent are fixed
before the centre accuracy is chosen. With
$N=4M+40$, the residual, corrected derivative and second derivative
satisfy the three gates of Theorem~\ref{inf:newton-threshold} on a
ball of radius $|\theta|p^{-2M-62/3}$.
The exact zero remains strictly convex and keeps a nonzero
physical monic seed coefficient. Its coefficient equation is
identified with a weak equation on the completed spaces, then
with a classical solution by elliptic regularity.

The logical order of the main steps is
\[
 \begin{gathered}
 \text{uniform Bessel expansion and congruence approximation}
 \ \Longrightarrow\ \mathcal P_{\mathrm{bd}},\ R,\ \theta,\\
 \text{angular positivity and exact local spectral flow}
 \ \Longrightarrow\ \text{integer-split gap},\\
 \text{finite centres and exact clamping}
 \ \Longrightarrow\ \text{residuals and material map},\\
 \text{relative graph bounds and bounded-cardinality clusters}
 \ \Longrightarrow\ \text{coefficient inverse},\\
 \text{Newton threshold}
 \ \Longrightarrow\ \text{classical convex non-ball domain}.
 \end{gathered}
\]
For Theorem~\ref{inf:theorem-B}, one radial pole per angular degree
is retained first. A finite-shift Lommel recurrence isolates four
critical phase ratios. Their short components have subpolynomial
coefficient-to-Hilbert conditioning, while the complementary
couplings have two inverse powers of $p$, up to a chosen small
loss. The exact centre is restored in the mixed norm before the
single radial conversion factor $\sqrt p$.
Section~\ref{inf:sec-sharp} gives material and nonlinear exponents
two and three; Section~\ref{inf:sec-tame} records their dependence
on $\tau\le\log p$. The three gates then close at order twenty
and radius $\tau p^{-11}$ as in
Theorem~\ref{inf:fixed-order-newton}.

The notation index in Appendix~\ref{inf:sec-notation} distinguishes
the fixed flow window, the local-count window, and the shrinking
inverse window.

\subsection{Classical inputs}

The proof uses the following classical results; the specialized
estimates needed in the construction
are proved in the body of the paper.
\begin{enumerate}[label=\textup{(\arabic*)}]
\item The integer-order common-zero theorem of Bourget, as a
consequence of Siegel's theorem, excludes integer primary
crossings \cite[\S15.28]{inf:watson}; see also
\cite{inf:siegel}. At half-integer orders the sine--cosine
representation and Hermite--Lindemann theorem
\cite[\S1.2]{inf:BakerWustholz} give non-coincidence as proved in
Lemma~\ref{inf:half-integer-noncoincidence}.
The finite Lommel polynomials \cite{inf:watson} give the
shift resonance estimate in Lemma~\ref{inf:lommel-resonance}.
The transcendence of $\pi$
\cite{inf:lindemann} makes $\alpha_{\mathrm{ar}}$ irrational.
\item Bessel power series and recurrences, the fixed-order outgoing Hankel
normalization and Watson's integral for the derivative of a
positive zero with respect to order are used in their classical
forms \cite[\S\S10.2, 10.6, 10.17; Eq.~10.21.17]{inf:dlmf}.
The uniform large-order remainders, derivative bounds, spacing,
and radial normalization used here are then proved in
Section~\ref{inf:sec-bessel}.
\item Gasper's positive linearization theorem is used with
\eqref{inf:gasper-hypothesis} and normalization at the endpoint of
the larger Jacobi exponent \cite[Theorem~1]{inf:gasper}.
The parameter substitution is given before
Lemma~\ref{inf:radial-stencils}.
\item Closed Dirichlet forms, relatively bounded perturbations,
compact resolvents, the min--max principle and finite Riesz
reductions are used as in \cite[Chs.~II and IV--VII]{inf:kato}.
The common domains, norm differentiability, eigenvalue derivative
argument and Schur reconstructions for this family are written
out in Section~\ref{inf:sec-flow}.
\item Dirichlet $H^2$ and higher regularity and Sobolev traces are
used on smooth bounded domains \cite[Chs.~6--9]{inf:GT}; analytic
boundary regularity is used for analytic boundaries and constant
Dirichlet data \cite{inf:MorreyNirenberg,inf:MorreyBoundary}.
The convex Dirichlet Hessian inequality is proved by its boundary
identity in \eqref{inf:convex-Hessian}. Polynomial core density
and the zero-capacity cutoff statements at the axes and origin
are included in Proposition~\ref{inf:flow-form}.
\item Hadamard's convexity theorem for a compact hypersurface
with positive second fundamental form is used in
Lemma~\ref{inf:geometry-convexity}; see
\cite{inf:doCarmoLima,inf:sacksteder}.
Bernstein's inequality for trigonometric polynomials, applied on
great circles, gives the dimension-independent shape derivative
bounds \cite[Vol.~II, Ch.~X, Theorem~3.13]{inf:zygmund}. The Jacobi endpoint bound itself is
proved in Lemma~\ref{inf:angular-sup}.
The rotation argument in Proposition~\ref{inf:arithmetic-density}
uses trigonometric polynomial approximation \cite{inf:zygmund}.
\item Banach's contraction theorem \cite{inf:banach} is applied to
a closed ball in the stated complete spaces.
Taylor's formula, continued-fraction determinant identities and
elementary lattice determinant facts are used in the explicit
proofs of their specialized forms.
\item The analytic implicit function theorem is used for the
scaled radius and seed equations in Theorem~\ref{inf:centres}.
Cauchy's estimates and the Hermite--Genocchi integral formula for
divided differences give the analytic coefficient bounds in
Lemma~\ref{inf:analytic-coefficients}; Cauchy's estimates also
give the centre remainder bounds.
The Dirichlet--Poincar\'e (Wirtinger) inequality on an interval is
used in Lemma~\ref{inf:bessel-spacing} and
Theorem~\ref{inf:local-count}.
\end{enumerate}
